\documentclass[11pt]{article}

\usepackage[margin=1in]{geometry}
\usepackage{amsmath}
\usepackage{amssymb}
\usepackage[T1]{fontenc}
\usepackage[utf8]{inputenc}
\usepackage{lmodern}

\title{Thermal Hydraulics for Reactor Kinetics Lab}
\author{}
\date{}

\begin{document}

\maketitle

\section{Purpose}

This note is a first-pass thermal-hydraulics design reference for the
\emph{Reactor Kinetics Lab} repository.  Its purpose is not to produce a final
plant design.  Instead, it does three narrower things:

\begin{enumerate}
  \item explains why the current Modelica prototype overheats so quickly,
  \item derives rough coolant-capacity targets from the repository's existing
        20 MW annular-core assumptions,
  \item converts those targets into a more realistic Modelica subsystem
        structure for the next implementation pass.
\end{enumerate}

The intended use is practical: this note should give enough engineering scale to
replace the current toy loop with a research-reactor-like thermal-hydraulic
model that stays within plausible operating temperatures.

\section{Current mismatch}

The current Python neutronics model assumes a reactor with:

\begin{itemize}
  \item nominal thermal power of 20 MW,
  \item annular active core geometry with $R_{\mathrm{inner}} = 0.8$ m,
        $R_{\mathrm{fuel}} = 3.456$ m, and active height $H = 6.912$ m,
  \item heavy-water moderation as the frozen physics assumption.
\end{itemize}

The current Modelica primary-loop prototype is useful as a numerical proof of
concept, but it is still a small pedagogical loop.  In particular, the present
prototype:

\begin{itemize}
  \item uses a single equivalent heated pipe rather than a reactor-scale coolant
        channel network,
  \item uses a simple radiator-like sink rather than a reactor heat exchanger,
  \item has no explicit saturation-margin target,
  \item does not yet enforce a reactor-appropriate pressure regime.
\end{itemize}

That is why the hot-leg temperature can shoot toward values such as
$400^{\circ}\mathrm{C}$ even while the return side remains much cooler.  The
prototype is absorbing heat numerically, but it is not yet constrained by the
real coolant inventory, realistic flow capacity, or a physically chosen primary
operating envelope.  A real low-pressure research reactor would not tolerate a
bulk outlet temperature anywhere near this regime; long before that point,
boiling, voiding, structural damage, and severe operational upset would occur.

\section{Reactor-scale inputs}

For this note, the frozen reactor-scale inputs are taken directly from the
existing repository documentation:

\begin{center}
\begin{tabular}{lc}
\hline
Quantity & Value \\
\hline
Nominal thermal power $P$ & $20$ MW \\
Inner moderator radius $R_{\mathrm{inner}}$ & $0.8$ m \\
Fuel outer radius $R_{\mathrm{fuel}}$ & $3.456$ m \\
Active height $H$ & $6.912$ m \\
Primary coolant assumption & D$_2$O \\
\hline
\end{tabular}
\end{center}

The active annular volume is approximated by
\[
V_{\mathrm{core}} = \pi \left(R_{\mathrm{fuel}}^2 - R_{\mathrm{inner}}^2\right) H.
\]

Using the current geometry,
\[
V_{\mathrm{core}} \approx \pi \left(3.456^2 - 0.8^2\right) \times 6.912
\approx 245.5\ \mathrm{m^3}.
\]

This gives an average volumetric power density of
\[
q''' = \frac{P}{V_{\mathrm{core}}}
\approx \frac{20 \times 10^6}{245.5}
\approx 8.15 \times 10^4\ \mathrm{W/m^3}
= 0.0815\ \mathrm{MW/m^3}.
\]

This value is not extreme for a large, heavily moderated core.  The main issue
is therefore not that the core is too power-dense, but that the present
prototype loop is not sized to move reactor-scale heat with reactor-scale
coolant flow.

\section{Cooling-capacity estimates}

\subsection{Bulk coolant temperature rise}

For a first-pass heavy-water estimate, it is acceptable to use water-like bulk
properties for sizing:

\[
c_p \approx 4.2\ \mathrm{kJ/(kg\,K)},
\qquad
\rho \approx 1100\ \mathrm{kg/m^3}.
\]

The required mass flow for a chosen bulk temperature rise is approximated by
\[
\dot{m} = \frac{Q}{c_p \Delta T}.
\]

For $Q = 20$ MW, the required total primary-loop flow is:

\begin{center}
\begin{tabular}{lcc}
\hline
Bulk rise $\Delta T$ & $\dot{m}$ [kg/s] & $\dot{V}$ [m$^3$/s] \\
\hline
10 K & 476.2 & 0.4329 \\
20 K & 238.1 & 0.2165 \\
30 K & 158.7 & 0.1443 \\
40 K & 119.0 & 0.1082 \\
50 K & 95.2 & 0.0866 \\
\hline
\end{tabular}
\end{center}

The most useful design range for the next Modelica pass is the middle of this
table, not the extremes.  A 20 K to 30 K bulk rise is a reasonable first target
for a research-reactor-like forced-circulation system.  That implies a total
primary mass flow on the order of 160 to 240 kg/s.

\subsection{Recommended bulk temperature envelope}

Because the current repository concept is much closer to a research reactor than
to a high-temperature power reactor, the next Modelica model should target a
bulk operating envelope well below the saturation concerns implied by a
$400^{\circ}\mathrm{C}$ hot leg.  A practical first-pass envelope is:

\begin{center}
\begin{tabular}{lc}
\hline
Quantity & Suggested target \\
\hline
Primary inlet bulk temperature & $45$--$60^{\circ}\mathrm{C}$ \\
Primary outlet bulk temperature & $65$--$90^{\circ}\mathrm{C}$ \\
Bulk rise across core & $20$--$30$ K \\
Pressure regime & low-pressure or modestly pressurized \\
\hline
\end{tabular}
\end{center}

This does not prove safety margins, but it does set a sane target for the next
Modelica architecture: the hot leg should stay in a low-temperature,
single-phase, research-reactor regime rather than drifting toward a PWR-like
high-temperature regime that the current design does not support.

\subsection{Heat-transfer surface argument}

If the annulus were naively cooled only through its inner and outer cylindrical
boundaries, the corresponding bounding area would be
\[
A_{\mathrm{bounds}} = 2\pi (R_{\mathrm{inner}} + R_{\mathrm{fuel}}) H
\approx 184.8\ \mathrm{m^2}.
\]

At 20 MW, that would correspond to an average heat flux of
\[
q'' \approx \frac{20 \times 10^6}{184.8} \approx 108\ \mathrm{kW/m^2}.
\]

This is already enough to show that the next Modelica model should not be built
as one tiny equivalent pipe with one tiny equivalent wall area.  A real core
would expose a much larger wetted area through many coolant channels, plates, or
subassemblies.  Therefore the next thermal-hydraulic abstraction should move
toward \emph{parallel channels plus plena}, not a single narrow toy conduit.

\subsection{Pump-duty scale}

The hydraulic power scale is also useful for planning.  Using
\[
W_{\mathrm{pump}} \approx \frac{\Delta p\,\dot{V}}{\eta},
\]
with $\eta \approx 0.7$ and the flow rates above, the shaft-power scale is on
the order of:

\begin{center}
\begin{tabular}{lccc}
\hline
Bulk rise $\Delta T$ & $\Delta p = 0.5$ MPa & $1.0$ MPa & $2.0$ MPa \\
\hline
20 K & 155 kW & 309 kW & 618 kW \\
30 K & 103 kW & 206 kW & 412 kW \\
40 K & 77 kW & 155 kW & 309 kW \\
\hline
\end{tabular}
\end{center}

These are only order-of-magnitude values, but they are already enough to guide
the next Modelica system away from a tiny pump and toward a reactor-scale head
curve with meaningful coastdown behavior.

\section{Candidate primary-loop architectures}

\subsection{Option A: low-pressure research-reactor concept}

The first candidate is a low-pressure, large-inventory heavy-water system with
forced circulation through a channelized core and a separate heat sink.  This is
the option that best matches the current repository's research-reactor framing.

Its main features would be:

\begin{itemize}
  \item inlet and outlet plena with significant coolant inventory,
  \item many effective coolant channels through the core,
  \item a forced-circulation pump train,
  \item a primary heat exchanger or cooler outside the core,
  \item bulk temperatures kept comfortably below boiling.
\end{itemize}

This option is the most natural follow-on to the existing open-tank intuition,
but it must become much more structured than the current prototype.  In
particular, the ``open tank + radiator'' picture should evolve into
``plenum/vessel volumes + channel bank + dedicated heat sink.''

\subsection{Option B: closed forced-circulation loop with secondary side}

The second candidate is a more explicitly closed D$_2$O primary loop feeding a
secondary heat exchanger.  This offers cleaner control over pressure and heat
rejection, and it is easier to connect to a later FMU co-simulation because the
boundaries are more explicit.

Its main features would be:

\begin{itemize}
  \item a closed primary inventory,
  \item one or more pumps with a realistic head curve,
  \item explicit primary-side heat exchanger,
  \item secondary-side sink represented either by a fixed-temperature sink or a
        simplified secondary loop.
\end{itemize}

This option is more structured and may eventually be better for FMI coupling,
but it also moves the model farther from the very simple research-reactor-like
mental model the current prototype was trying to preserve.

\section{Recommended operating envelope and structure}

For the next Modelica implementation, the recommended direction is:

\begin{quote}
Use a \textbf{low-temperature forced-circulation D$_2$O primary model} with
\textbf{channelized core flow}, \textbf{plena or vessel volumes}, and an
\textbf{explicit heat sink}, while keeping bulk outlet temperatures in the
roughly $65$--$90^{\circ}\mathrm{C}$ range rather than allowing a
high-temperature pressurized-loop regime to emerge by accident.
\end{quote}

That recommendation is based on three observations:

\begin{enumerate}
  \item the reactor-scale power density is modest enough that a large-flow,
        low-temperature solution is plausible,
  \item the current repository already assumes a research reactor, not a compact
        power reactor,
  \item the present Modelica overheating is better interpreted as a missing
        thermal-hydraulic capacity model than as evidence that the reactor must
        operate at very high primary temperatures.
\end{enumerate}

\section{Modelica implementation targets}

The next thermal-hydraulic Modelica model should no longer be a single heated
pipe plus radiator.  A better first abstraction is:

\begin{enumerate}
  \item \textbf{inlet plenum or lower vessel volume},
  \item \textbf{equivalent channel bank} through the core, either as parallel
        channels or one channel with explicit scaling parameters,
  \item \textbf{outlet plenum or upper vessel volume},
  \item \textbf{pump model} with a head curve and optional coastdown,
  \item \textbf{primary heat exchanger or cooler} that removes the full 20 MW,
  \item \textbf{pressure boundary strategy} consistent with a low-pressure or
        modestly pressurized research-reactor system.
\end{enumerate}

The minimum FMI-oriented variable set should be:

\begin{center}
\begin{tabular}{ll}
\hline
Direction & Variables \\
\hline
Python $\rightarrow$ Modelica & total core power, axial power fractions, control actions \\
Modelica $\rightarrow$ Python & inlet temperature, outlet temperature, total flow, pressure drop \\
\hline
\end{tabular}
\end{center}

For the channel abstraction, the following scaling logic is recommended:

\begin{itemize}
  \item total loop flow should be a reactor-scale quantity, about 160--240 kg/s
        for a 20--30 K rise,
  \item core heating should be distributed axially, as already planned,
  \item the effective coolant area and wetted surface should represent many
        channels, not one narrow laboratory pipe.
\end{itemize}

\section{Immediate implications for the current prototype}

The present Modelica prototype is still useful, but it should now be interpreted
as a \emph{numerical topology prototype}, not as a thermal-hydraulic design.
Before adding full neutronics feedback, the next Modelica pass should at least:

\begin{enumerate}
  \item scale the nominal power from 10 kW-like behavior to the full 20 MW,
  \item replace the tiny equivalent core pipe with a channel-bank abstraction,
  \item replace the radiator sink with a reactor-scale heat sink or exchanger,
  \item enforce a low-temperature operating envelope that prevents unrealistic
        hot-leg temperatures.
\end{enumerate}

\section{Open questions}

This note deliberately leaves several questions open for the next iteration:

\begin{itemize}
  \item Should the next primary model be low-pressure and open/pool-like, or
        slightly pressurized but still firmly in a research-reactor regime?
  \item How many effective parallel coolant channels are needed to make the
        channel-bank abstraction believable without introducing unnecessary
        numerical stiffness?
  \item Is a heavy-water property package needed immediately, or is water-like
        property approximation acceptable for the first reactor-scale thermal
        sizing pass?
\end{itemize}

The key point is that these are now engineering-structure questions, not a sign
that the reactor should be allowed to run at several hundred degrees Celsius in
the primary loop.

\end{document}